\section{Solutions for Chapter~\ref{ch:factors}}

\paragraph{Exercise~\ref{ex:c-factors-1}: parafermion charges.}
\phantomsection
\label{sol:ex:c-factors-1}
Subtracting the charge-one Cartan stress gives
$c_{\mathrm{PF}_k}=3k/(k+2)-1=2(k-1)/(k+2)$.
Thus the three substitutions are $2/4=1/2$,
$4/5$, and $24/15=8/5$, respectively.
The first is the Ising parafermion algebra. Notice that the
parafermion central charges approach two from below as $k$ tends
to infinity; proximity to two by itself therefore selects no
particular level or full theory.



\paragraph{Exercise~\ref{ex:c-factors-2}: three opposite momentum pairs.}
\phantomsection
\label{sol:ex:c-factors-2}
Using $r^2=13$, $t^2=-3$, and $\overline t=-t$, the three products are
\[
 p_{\beta_1}p_{-\beta_1}
 =p_{\beta_3}p_{-\beta_3}
 =\frac{r^2-9t^2}{1000}
 =\frac{40}{1000}=\frac1{25},
\]
and
\[
 p_{\beta_2}p_{-\beta_2}
 =\frac{49-9r^2t^2}{10000}
 =\frac{49+351}{10000}=\frac1{25}.
\]
The transported PCT multiplies an opposite exponential by
$10,1,10$ for the three positive momenta. Directly,
\begin{align*}
 10\overline{p_{\beta_1}}&=(-r+3t)/10=p_{-\beta_1},\\
 \overline{p_{\beta_2}}&=(7-3rt)/100=p_{-\beta_2},\\
 10\overline{p_{\beta_3}}&=-(r+3t)/10=p_{-\beta_3}.
\end{align*}
The opposite directions use the reciprocal factors, so each
two-term exponential pair is fixed. Each Heisenberg oscillator
changes sign under PCT; its quadratic products are fixed as well.
Consequently the entire Potts stress is PCT real in the transported
positive form.



\paragraph{Exercise~\ref{ex:c-factors-3}: the spin-three simple pole.}
\phantomsection
\label{sol:ex:c-factors-3}
Skew symmetry gives
$a_{(0)}a=-a_{(0)}a+D(a_{(1)}a)-\frac12D^2(a_{(2)}a)
+\frac16D^3(a_{(3)}a)-\cdots$ for an even field $a$.
Here ``even'' means bosonic parity, despite the odd conformal
weight three; no fermionic minus sign is present.
Use $a_{(1)}a=b_4$, $a_{(2)}a=DL$, $a_{(3)}a=2L$.
Derivatives of the scalar higher pole vanish. Thus
\[
 2a_{(0)}a=Db_4-\frac12D^3L+\frac13D^3L
            =Db_4-\frac16D^3L.
\]
Dividing by two gives the displayed self-row.



\paragraph{Exercise~\ref{ex:c-factors-4}: the two isotropic lines.}
\phantomsection
\label{sol:ex:c-factors-4}
If $j=0$, isotropy forces $k=0$. Otherwise divide by $j^2$ to
obtain $(k/j)^2=-1$ in $\mathbb F_{13}$. The two roots are
$5$ and $-5$. Adjoining both lines would adjoin their span, which
is the entire two-dimensional discriminant group. It contains
$(1,0)$, whose quadratic value is $-1/13$, so such an extension
would not have integral vacuum grading.
Explicitly $5^2=25\equiv-1\pmod{13}$, and
$(1,5),(1,-5)$ are independent because their determinant
$-10$ is nonzero modulo thirteen. A quadratic polynomial over a
field has at most two roots, so there are no further isotropic lines.



\paragraph{Exercise~\ref{ex:c-factors-5}: coefficients of the stress identity.}
\phantomsection
\label{sol:ex:c-factors-5}
In the basis $(\omega_L,\omega_R,e,f)$, the four coefficients on
the proposed right side are
\begin{align*}
 \omega_L &: 1,\\
 \omega_R &: \frac{3403}{4563}+\frac5{27}+\frac{35}{507}
             =\frac{3403+845+315}{4563}=1,\\
 e &: \frac38+\frac58=1,\\
 f &: \frac{13}{14}+\frac1{14}=1.
\end{align*}
They therefore reproduce the sum of the four orthogonal stresses.
The norm of a vector gives only one quadratic constraint; it
does not determine these four projection coefficients. More
fundamentally, a correct central term is only one pole in a
self-OPE. Establishing the Potts or parafermion subalgebra also
requires the remaining singular products, the requisite higher
fields, and the stated translation and PCT properties.



\paragraph{Exercise~\ref{ex:c-factors-6}: positivity without commutativity.}
\phantomsection
\label{sol:ex:c-factors-6}
Take
\[
 A=\begin{pmatrix}2&0\\0&1\end{pmatrix},
 \qquad B=\begin{pmatrix}2&1\\1&2\end{pmatrix}.
\]
Their spectra are $\{1,2\}$ and $\{1,3\}$, respectively, so both
are positive definite and self-adjoint. But
$[A,B]=\left(\begin{smallmatrix}0&1\\-1&0\end{smallmatrix}\right)$.
For $\alpha,\beta>0$ and every vector $v$,
\[
 (\alpha+\beta)\|v\|^2
 \leq
 \langle(\alpha A+\beta B)v,v\rangle
 \leq(2\alpha+3\beta)\|v\|^2.
\]
This follows by bounding each expectation separately and adding.
No common eigenbasis was used. The $C_2$ proof does the same
thing with the partial energies on a fixed ambient grade: it bounds
each expectation on the same vector in $C_2(V)^\perp$, then uses
the positive stress identity to bound the total energy.


